\begin{table}[htbp]\centering\small\setlength{\tabcolsep}{4pt}
\caption{两组与三组方案中各执行组的独立资源配置}\label{tab:q4_each_group}
\begin{tabular}{llrrrrrrrrr}\toprule
分组 & 服务区 & \multicolumn{3}{c}{运输机/架} & \multicolumn{3}{c}{电池/组} & 中继机 & 组件 & 作业/min\\
 & & A & B & C & A & B & C & /架 & /组 & \\\midrule
2组-G1 & 除S006外14区 & 4 & 2 & 2 & 6 & 4 & 3 & 2 & 3 & 867.745\\
2组-G2 & S006 & 0 & 0 & 1 & 0 & 0 & 1 & 0 & 0 & 26.026\\
3组-G1 & S001 & 1 & 0 & 2 & 1 & 0 & 2 & 0 & 0 & 71.754\\
3组-G2 & 其余13区 & 4 & 2 & 2 & 6 & 4 & 3 & 2 & 3 & 795.991\\
3组-G3 & S006 & 0 & 0 & 1 & 0 & 0 & 1 & 0 & 0 & 26.026\\
\bottomrule\end{tabular}
\par\smallskip\footnotesize 两组与三组的G编号分别定义，不是跨方案固定实体；作业时间不包含充电和返航周转，资源配置区间则按各类实际释放规则核算。
\end{table}
